\documentclass[10pt,a4paper]{amsart}
\usepackage[margin=19mm]{geometry}
\usepackage{amsmath,amssymb,mathtools}
\usepackage[scr=rsfs,cal=euler]{mathalfa}
\usepackage{xfrac}
\usepackage[unicode,hidelinks,bookmarks=false]{hyperref}
\newcommand{\ie}{i.e.\ }
\newcommand{\cf}{cf.\ }
\newcommand{\resp}{respectively\ }
\newcommand{\Subsection}{\subsection}
\providecommand{\nobreakdash}{\nobreak\hbox{-}}
\setlength{\emergencystretch}{2em}
\multlinegap0pt
\usepackage{bm}
\usepackage{bbm}
\usepackage{dsfont}
\usepackage{xspace}
\usepackage{cancel}
\usepackage{mathtools}
\usepackage{amsthm}
\makeatletter
\@ifundefined{theorem}{\newtheorem{theorem}{Theorem}}{}
\makeatother
\renewcommand{\qedsymbol}{}  % Removes the box for all proofs
\title[On the Intractability of Chaotic Symbolic Walks: Toward a No]{On the Intractability of Chaotic Symbolic Walks: Toward a Non-Algebraic Post-Quantum Hardness Assumption}
\author{Mohamed Aly Bouke}
\date{Source: arXiv:2505.22644v1 (CC BY 4.0)}
\begin{document}
\maketitle

\noindent\emph{Source and licence.} On the Intractability of Chaotic Symbolic Walks: Toward a Non-Algebraic Post-Quantum Hardness Assumption. Version arXiv:2505.22644v1 (CC BY 4.0), distributed under the Creative Commons Attribution 4.0 International licence, as recorded in the arXiv licence metadata for this exact version and shown on the abstract page https://arxiv.org/abs/2505.22644v1. Full rights notice: licenses/arxiv-2505.22644-CC-BY-4.0.txt.\par\smallskip

\noindent\textit{Editorial scope.} Counted unit: the Theorem whose source label is \texttt{None} in the article (source0.tex lines 464--466, source label \texttt{None}), delivered together with the article's own complete original proof of it (source0.tex lines 468--491, one \texttt{proof} environment of 24 lines). No mathematics is written, repaired or re-derived: statement and proof are the article's own text line for line. Required context retained uncounted: none. Every definition, display and label of those lines is retained unchanged. Omitted from this package and not redistributed: 1--463, 467--467, 492--809. Nothing inside a retained excerpt is omitted or altered: every retained body is delivered exactly as the article's lines are written, so no omission receipt of any kind accompanies this package. Citations: the retained excerpts carry no citation command at all, so no bibliography entry of the article is transcribed and no reference is redistributed. Reference adaptation: none was needed and none was made; every reference inside the delivered unit resolves inside it, so no unresolved reference or citation is produced. Printed slips kept literal: no printed word, symbol, hypothesis or number of the counted statement or of its proof was altered. Layout and class: the article's own class is \texttt{article} with packages including fontenc, inputenc, mathptmx, helvet, courier, geometry, graphicx, xcolor, cuted, microtype, parskip, tcolorbox, academicons, hyperref; the wrapper is the lane's pinned \texttt{amsart} editorial wrapper at 10pt on A4, which restates the theorem-like environments the retained text needs and adds the article's own macro definitions that the retained text needs (\texttt{\textbackslash qedsymbol}), with extra wrapper packages (bm, bbm, dsfont, xspace, cancel, mathtools). The delivered numbering is the unit's own. Assets: no retained span contains a figure environment, a graphics-inclusion command, a PDF-page inclusion, a table or a TikZ picture; no image file is copied into this package, the package declares no asset dependency and no image file is redistributed with it. Licence: version arXiv:2505.22644v1 of this article is distributed under the Creative Commons Attribution 4.0 International licence, as recorded in the arXiv licence metadata for this exact version and shown on the abstract page https://arxiv.org/abs/2505.22644v1. A full rights notice is delivered at licenses/arxiv-2505.22644-CC-BY-4.0.txt. Characters: every retained excerpt is pure ASCII, so the delivered engine, XeLaTeX with its default Latin Modern text font, reports no missing character.
\begin{theorem}
The problem of counting valid SPIP trajectories (SPIP-Count) is \#P-hard.
\end{theorem}

\begin{proof}
Let $G = (V, E)$ be an instance of \#PATH COUNT. Our goal is to construct an SPIP instance in which the number of valid symbolic trajectories from $x_0$ to $x_n$ corresponds exactly to the number of $s \rightarrow t$ paths in $G$.

\textbf{Step 1: State Encoding.} Assign each vertex $v_i \in V$ a unique symbolic state $x_i \in \mathbb{Z}^2$ via a bijection $\phi: V \rightarrow \mathbb{Z}^2$.

\textbf{Step 2: Edge Encoding.} For each edge $(u,v) \in E$, define a contractive affine map $T_{(u,v)}$ such that:
\[
T_{(u,v)}(x_u) = \lfloor A x_u + b_{(u,v)} + \delta_{(u,v)} \rfloor \approx x_v
\]
Where:
\begin{itemize}
    \item $A$ is a fixed contractive matrix (e.g., $0.5I$),
    \item $b_{(u,v)}$ is chosen so that the unperturbed map lands near $x_v$,
    \item $\delta_{(u,v)}$ is bounded noise (e.g., in $[-1,1]^2$) allowing rounding to $x_v$.
\end{itemize}

\textbf{Step 3: Trajectory Construction.} A valid symbolic trajectory in SPIP corresponds to a valid path in $G$. Each sequence of transformations that maps $x_0 = \phi(s)$ to $x_n = \phi(t)$ represents a symbolic realization of a path from $s$ to $t$.

\textbf{Step 4: One-to-One Mapping.} For each path $\pi = (s \rightarrow v_1 \rightarrow \dots \rightarrow t)$ in $G$, there exists a unique (up to bounded noise) symbolic trajectory $\mathcal{P}_\pi$ in SPIP that applies the corresponding affine maps.

\textbf{Step 5: Reduction.} Since each valid path in $G$ maps to a distinct trajectory in SPIP, counting SPIP trajectories from $x_0$ to $x_n$ is equivalent to counting $s \rightarrow t$ paths in $G$. The reduction can be computed in polynomial time.

Hence, SPIP-Count is \#P-hard. \renewcommand{\qedsymbol}{}
\end{proof}

\end{document}
